% =============================================================================
% Bloque 1: Apertura (slides 1–2)
% =============================================================================

% --- Slide 1: Portada PhytoLearning (charla de divulgación) ---
\begin{frame}[plain,noframenumbering]
\vspace*{0.6em}
\begin{columns}[T,onlytextwidth]
%
\begin{column}{0.68\textwidth}
\raggedright
{\color{uaigray}\footnotesize Reunión Mensual \textbf{PhytoLearning}}

\vspace{1.1em}

{\color{uaibluedark}\fontsize{20}{24}\selectfont\bfseries Geometría y LLMs\\ para clases con pocos ejemplos}\\[0.55em]
{\color{uaigray}\small Cómo elegir datos generados por IA}

\vspace{1.0em}
{\color{uaiorange}\rule{0.52\textwidth}{1pt}}
\vspace{1.0em}

{\large\bfseries\color{uaibluedark} Benjamín Schindler}\\[4pt]
{\footnotesize\color{uaigray}
Laboratorio Dr.\ Gonzalo A.\ Ruz\\[2pt]
Master of Science in Data Science\\
Universidad Adolfo Ibáñez\par}
\end{column}
%
\begin{column}{0.30\textwidth}
\raggedleft
\includegraphics[height=1.0cm]{logos/fic.png}
\end{column}
\end{columns}
\end{frame}

% --- Slide 2: Roadmap (lista con descripciones, sin tiempos) ---
\begin{frame}{Hoja de ruta}
\small
\setbeamercolor{enumerate item}{fg=uaiorange}
\begin{enumerate}\setlength\itemsep{4pt}
  \item \textbf{Introducción} — desbalance, voces que se callan y por qué las clases raras suelen ser las que más importan.

  \item \textbf{Problema} — qué pasa cuando hay pocos ejemplos y qué se hace hoy al respecto.

  \item \textbf{Estado del arte} — dos caminos para inventar ejemplos (el clásico y la IA) y lo que le falta a ambos.

  \item \textbf{Hipótesis y objetivos} — qué afirmamos y cómo se va a verificar.

  \item \textbf{Metodología} — un «mapa de significados», filtros de calidad sobre ese mapa y pesos por calidad.

  \item \textbf{Resultados} — \alert{8.000+ experimentos}: la idea funciona, y lo simple le gana a lo complejo.

  \item \textbf{Discusión} — utilidad práctica, limitaciones honestas y trabajo futuro.
\end{enumerate}
\end{frame}
